\begin{abstract}
Financial multi-agent systems have achieved notable progress in role specialization and textual reasoning, yet their self-improvement mechanisms remain confined to verbal reinforcement, where reflections are rewritten into prompt prefixes without penetrating model parameters or consolidating into long-term memory. This limitation becomes critical under regime shifts: recent live benchmarks reveal that even frontier models incur net losses in real trading windows, and reinforcement-finetuned agents exhibit improved classification accuracy but degraded cumulative returns. We present \textsc{FinOPD}, a financial multi-agent framework whose core mechanism is experience-driven self-evolution through on-policy self-distillation. The same model serves as both student and teacher under differential conditioning: the student reasons from current observations while the teacher additionally conditions on hindsight risk-adjusted returns as privileged information. A token-level Jensen-Shannon divergence objective, combined with per-token KL clipping and Shapley-based credit assignment across five specialist agents, drives parametric evolution without requiring an external teacher. In parallel, a non-parametric track consolidates winning geometry-factor-action triples into a retrieval-augmented belief memory indexed by visual-geometric embeddings. Across a dual-market benchmark spanning Dow-30 equities and CSI-300 constituents, \textsc{FinOPD} demonstrates consistent improvements in Sharpe ratio, maximum drawdown, and cross-regime robustness over twelve baselines including TradingAgents, FinCon, and AlphaAgent. The evolution curve exhibits observable monotonic improvement across self-distillation iterations under three random seeds, a property that persists on a strictly post-cutoff forward evaluation window, ruling out pretraining leakage as an explanation.
\end{abstract}
